\documentclass[11pt,a4paper]{article}

\usepackage[utf8]{inputenc}
\usepackage[T1]{fontenc}
\usepackage[french]{babel}
\usepackage{amsmath}
\usepackage[margin=2.4cm]{geometry}
\usepackage{booktabs}
\usepackage{array}
\usepackage{ragged2e}
\usepackage{sectsty}
\usepackage{setspace}
\usepackage{microtype}
\usepackage[colorlinks=true,linkcolor=black,urlcolor=black]{hyperref}
\usepackage{titlesec}

\allsectionsfont{\bfseries}
\titleformat{\section}{\normalsize\bfseries}{\arabic{section}.}{0.5em}{}
\titleformat{\subsection}{\normalsize\bfseries}{\arabic{section}.\arabic{subsection}}{0.5em}{}
\setstretch{1.2}
\newcolumntype{L}[1]{>{\RaggedRight\arraybackslash}p{#1}}
\renewcommand{\arraystretch}{1.3}

\begin{document}

% ======================= ARTICLE =======================
\begin{center}
  {\LARGE\bfseries La completude logique des interviews COSO\par}
  \vspace{0.5em}
  {\normalsize Une mesure de la fiabilite de la collecte par reconstruction
            des questions attendues\par}
  \vspace{1.2em}
  {\normalsize \textbf{Note d'analyse} --- Enquete COSO, base
            \texttt{Questionnaire\_COSO\_VF.dta} (1032 interviews)\par}
\end{center}
\vspace{1.2em}

% ---------------- RESUME ----------------
\noindent\begin{minipage}{\textwidth}
\textbf{Resume.} Dans les enquetes a logique de saut, le taux de reponse
brut ne renseigne pas sur la qualite reelle de la collecte : une interview
``complete'' au sens administratif peut laisser vides des questions qui
devaient pourtant etre posees. Cette note reconstruit, pour chaque
interview, l'ensemble des questions attendues au regard des reponses
anterieures, puis mesure combien d'entre elles ont reellement recu une
reponse. La methode est validee sur un groupe de reference de
425 interviews controles par le superviseur. Les resultats indiquent que
\textbf{627 des 1032 interviews (60,8~\%) sont totalement remplies}, et que
rejoindre les quasi-completes porte la base fiable a 689 interviews.
La methode fournit aussi, pour chaque interview et chaque question, la
liste des manques, exploitable pour cibler des questionnaires a revoir.

\vspace{0.4em}
\textbf{Mots-cles :} completude ; logique de saut ; qualite de la collecte ;
enquete menages ; COSO.
\end{minipage}

\vspace{1.2em}
\hrule
\vspace{1.2em}

% ----------------- 1. INTRODUCTION -----------------
\section{Introduction et question}

Dans une enquete a questionnaires complexes, certaines questions ne sont
posees que lorsque la reponse a une question anterieure le justifie
(saut conditionnel). Le simple fait qu'une interview soit cataloguee comme
``conduite jusqu'au bout'' ne garantit pas que toutes les questions
obligatoires ont ete repondues : des abandons partiels, des refus locaux ou
des erreurs de saisie laissent des cellules vides la ou une reponse etait
attendue.

La question que se pose l'etude est donc : \emph{combien d'interviews sont,
sous la logique du questionnaire, totalement remplies ?} La reponse importe
pour la fiabilite des analyses : une base melee d'interviews completes et
incompletes biaise les statistiques par une absence non aleatoire de
reponses.

% ----------------- 2. DONNEES ET METHODE -----------------
\section{Donnees et methode}

\subsection{Donnees}

La base d'etude contient \textbf{1032 interviews}. Elle est issue de la
fusion des versions 4 et 5 du questionnaire et comporte 218 variables.
Deux elements de la collecte sont exploites pour la validation : le statut
de l'interview (\texttt{interview\_\_status}) et le resultat de l'appel
(A4), sans que ces deux variables n'entrent dans le calcul du taux de
completude lui-meme.

\subsection{Taux de completude}

Pour chaque interview $i$, on reconstruit l'ensemble $Q(i)$ des questions
attendues selon la logique de saut, puis on calcule :
%
\begin{center}
  $\displaystyle \text{taux}(i)=\frac{\text{nombre de questions attendues et
  repondues}}{|Q(i)|}$. \qquad (1)
\end{center}

Une reponse n'est comptee que si la cellule est non vide au sens strict :
les chaines vides \texttt{""} ou \texttt{.} --- presentes dans la base sans
etre des valeurs manquantes --- sont traitees comme des absences de
reponse. Les champs techniques et les champs libres hors logique (C6, C9,
C10, C11) sont exclus.

\subsection{Validation}

La reconstruction des regles comporte un risque d'erreur dans les deux
sens (regles trop strictes ou trop laxistes). Un \emph{metre etalon} est
donc constitue : les interviews validees par le superviseur et reellement
conduites (statut 100 et A4 = 1), soit 425 interviews. Ce groupe sert
uniquement a controler la methode : il est exclu du calcul des taux.
Si la reconstruction est juste, ces interviews doivent presque toutes
atteindre 100~\%.

% ----------------- 3. RESULTATS -----------------
\section{Resultats}

\subsection{Validite de la methode}

Le tableau~\ref{tab:valid} resume la validation sur les 425 references.
Le taux median de 1,000 et la proportion de 95,5~\% de references
exactement a 100~\% valident la reconstruction : la logique implantee n'est
ni trop stricte (sinon les references echoueraient), ni trop laxiste
(sinon elles seraient encore plus nombreuses a 100~\% sans que ce soit
justifie).

\begin{table}[h]
\centering
\caption{Validation sur le groupe de reference ($n=425$).}
\label{tab:valid}
\small
\begin{tabular}{L{8.4cm}c}
\toprule
\textbf{Indicateur} & \textbf{Valeur} \\
\midrule
References (statut 100 et A4 = 1)              & 425 \\
dont taux $= 100\%$                             & 406  (95,5~\%) \\
mediane du taux                                 & 1,000 \\
moyenne du taux                                 & 0,997 \\
references sous 90~\%                           & 2 \\
\bottomrule
\end{tabular}
\end{table}

\subsection{Repartition des interviews}

Le tableau~\ref{tab:classe} definit les cinq classes de completude ; le
tableau~\ref{tab:rep} en dresse le bilan sur les 1032 interviews.

\begin{table}[h]
\centering
\caption{Classes de completude et intervalles.}
\label{tab:classe}
\small
\begin{tabular}{L{3.4cm}L{4.0cm}L{5.2cm}}
\toprule
\textbf{Classe} & \textbf{Intervalle} & \textbf{Lecture} \\
\midrule
1-COMPLET        & taux $=100\%$               & aucune question attendue vide \\
2-QUASI\_COMPLET & $90\% \leq$ taux $<100\%$   & quelques manques residuels \\
3-MOYEN          & $50\% \leq$ taux $<90\%$    & une partie notable est vide \\
4-INCOMPLET      & $0\% <$ taux $<50\%$        & majorite des questions attendues vides \\
5-VIDE           & taux $=0\%$                 & aucune question exploitable \\
\bottomrule
\end{tabular}
\end{table}

\begin{table}[h]
\centering
\caption{Repartition des 1032 interviews.}
\label{tab:rep}
\small
\begin{tabular}{lcc}
\toprule
\textbf{Classe} & \textbf{Effectif} & \textbf{Part} \\
\midrule
1-COMPLET        & 627 & 60,8~\% \\
2-QUASI\_COMPLET & 62  & 6,0~\% \\
3-MOYEN          & 10  & 1,0~\% \\
4-INCOMPLET      & 291 & 28,2~\% \\
5-VIDE           & 42  & 4,1~\% \\
\midrule
TOTAL            & 1032 & 100,0~\% \\
\bottomrule
\end{tabular}
\end{table}

Le taux moyen est de 0,686 et la mediane de 1,000 : la distribution est
fortement biconnectee, avec une majorite d'interviews parfaitement
completes et une minorite quasi vides, refleterant des abandons precoces.

% ----------------- 4. DISCUSSION -----------------
\section{Discussion}

Le principal resultat est que \textbf{627 interviews (60,8~\%) sont
totalement remplies}, et que la base fiable (completes + quasi-completes)
reunit \textbf{689 interviews (66,7~\%)}. Ce bilan, plus exigeant que le
simple critere de consequence administrative, montre la valeur ajoutee
d'une mesure de completude : elle qualifie precisement ce qui est exploitable
au regard de la logique du questionnaire.

L'analyse identifie aussi, pour chaque interview, la liste exacte des
questions attendues et restees vides. Cette information permet de cibler,
par classe et par identifiant (\texttt{interview\_\_id} /
\texttt{interview\_\_key}), les questionnaires a revoir ou a recontacter,
et de reperer les questions les plus souvent laissees vides.

Le choix d'une mesure en flux (100~\% = complet === reste sont des
intervalles de tolerance) reste conventionnel : la classe
\texttt{2-QUASI\_COMPLET} (au moins 90~\%) offre un seuil pragmatique, car
certains champs << preciser >> facultatifs peuvent manquer sans que le
questionnaire soit reellement inacheve.

% ----------------- 5. CONCLUSION -----------------
\section{Conclusion}

La reconstruction de la logique conditionnelle du questionnaire COSO, puis
sa validation statistique sur 425 references controlees, permettent de
classifier les 1032 interviews en cinq niveaux de completude. La methode
repond a la question initiale --- 627 interviews totalement remplies --- et
fournit un outil operationnel (base filtrable et fichier Excel de detail)
pour la suite de l'exploitation des donnees.

\begin{thebibliography}{9}
\bibitem{coso}
  COSO. \emph{Questionnaire d'enquete --- Draft 8}, dossier
  \texttt{Outils\_de\_Collecte}, version 4-5 fusionnee.
\end{thebibliography}

\end{document}